\documentclass[10pt,a4paper]{amsart}
\usepackage[margin=19mm]{geometry}
\usepackage{amsmath,amssymb,mathtools}
\usepackage[scr=rsfs,cal=euler]{mathalfa}
\usepackage{xfrac}
\usepackage[unicode,hidelinks,bookmarks=false]{hyperref}
\newcommand{\ie}{i.e.\ }
\newcommand{\cf}{cf.\ }
\newcommand{\resp}{respectively\ }
\newcommand{\Subsection}{\subsection}
\providecommand{\nobreakdash}{\nobreak\hbox{-}}
\setlength{\emergencystretch}{2em}
\multlinegap0pt
\usepackage[all]{xy}
\usepackage[shortlabels]{enumitem}
\usepackage{xcolor}
\usepackage{amssymb}
\newtheorem{lemma}{Lemma}
\newtheorem{theorem}{Theorem}
\newcommand{\DL}{\mathcal{DL}}
\DeclareMathOperator{\Zinb}{{\mathrm{DL}}}
\newcommand{\cochain}[2]{C^{#1}_{\DL}(#2)}
\title{A New Approach to Defining Cochain Complexes for dual Leibniz algebra}
\author{Hassan Alhussein}
\date{Source: arXiv:2512.19319v1 (CC BY 4.0)}
\begin{document}
\maketitle

\noindent\emph{Source and licence.} A New Approach to Defining Cochain Complexes for dual Leibniz algebra. Version arXiv:2512.19319v1 (CC BY 4.0), distributed by arXiv under the Creative Commons Attribution 4.0 International licence, http://creativecommons.org/licenses/by/4.0/, as recorded in the arXiv licence metadata for this exact version and shown on the abstract page https://arxiv.org/abs/2512.19319v1. Full licence notice: \texttt{licenses/arxiv-2512.19319-CC-BY-4.0.txt}.\par\smallskip

\noindent\textit{Editorial scope.} Counted unit: the article's theorem 5.4 (source0.tex lines 324--330). The article's complete original proof of it is delivered with it (source0.tex lines 331--724, one \texttt{proof} environment of 394 lines). No mathematics is written, repaired or re-derived: the statement and the proof are the article's own lines, in the article's own notation, and no retained line is substituted or re-flowed.  Retained uncounted as the source-local context the proof needs: lines 115--128.  Nothing was omitted from any retained span, so the package carries no omission record.  The retained lines cite 1 works, whose entries are transcribed verbatim from the article's own bibliography source in the e-print and delivered with short editorial labels recorded in the item's reference mapping.  No retained line contains a figure environment, an image-inclusion command or drawing code, so no image is delivered and the package declares no asset dependency.  Editorial formatting only, in addition: an AMS \texttt{amsart} preamble that restates the article's own declarations and macros used by the retained lines, namely the environments \texttt{lemma}, \texttt{theorem} on separate counters, together with the standard packages those lines need.  The retained environments print in this document as Lemma 1, Theorem 1 in order of appearance, whereas the article prints the same items with the numbering of its own full document.  The complete native source of the article as distributed in the arXiv e-print for this exact version is bundled unchanged as \texttt{sources/191464/source0.tex}.
\begin{lemma}[\cite{I.S}]\label{lem}
  Let \(\mathfrak{g}\) be a Leibniz algebra and \(x, x_1, \dots, x_n \in \mathfrak{x}\). Then

\[
[x, [x_1, \dots, x_n]] = \sum_{i=0}^{n-1} \sum_{(\alpha,\beta) \in \mathsf{Sh}^1(n-i,i)} (-1)^i [x, x_{\beta(i)}, \dots, x_{\beta(1)}, x_{\alpha(1)}, \dots, x_{\alpha(n-i)}],
\]


Here \(\mathsf{Sh}^1(s,t)\) means \((s,t)\)-shuffles  and the left normed Leibniz  bracket of elements $x_1,\cdots,x_n$ of a Leibniz algebra  \(\mathfrak{g}\) is defined by recursion
\[
[x_1,\cdots,x_n] := [[x_1,\cdots,x_{n-1}], x_n],\quad where\quad [x_1] =x_1.
\]
  
\end{lemma}

\begin{theorem}\label{5.4}
Let $\mathfrak{g}$ be a Leibniz algebra, let \((B, \cdot_B)\) be a dual Leibniz algebra, and let \(M\) be a \(B\)-bimodule. Then there exists a cochain map
\[
\Psi : \cochain{n}{B, M} \rightarrow C^*_{\mathrm{Lie}}(\mathfrak{g} \otimes B, \mathfrak{g} \otimes M)
\]
from the dual Leibniz cochain complex of \(B\) with coefficients in \(M\) to the Lie cochain complex of the tensor product algebra  \(\mathfrak{g} \otimes B\) with coefficients in \(\mathfrak{g} \otimes M\).
\end{theorem}

\begin{proof}
  Given a dual Leibniz  cochain $f \in \cochain{n}{B, M}$, for $x_1,\ldots,x_n\in \mathfrak{g}$ and $b_1,\ldots,b_n\in B$, define: 
  \begin{align*}
\Psi(f)\big(x_1 \otimes b_1, x_2 \otimes b_2,\ldots,x_{n}\otimes b_{n}\big) : =\sum_{\sigma\in S_n}sgn(\sigma)[[[x_{\sigma(1)},x_{\sigma(2)}],\cdots,x_{\sigma(n-1)}],x_{\sigma(n)}]\otimes f(b_{\sigma(1)},\cdots,b_{\sigma(n)})
\end{align*}

Let us to show that $\Psi$ commutes with the  differentials:
\[
\delta_{Lie}(\Psi f)=\Psi(\delta_{\DL}f)
\]
For $n=2$ and $x =x_1  \otimes b_1$, $y = x_2 \otimes b_2$,
\[
\delta_{Lie}(\Psi f)(x,y) = \underbrace{x \ast (\Psi f)(y)}_{(1)} - \underbrace{(\Psi f)(x \ast y)}_{(2)}  - \underbrace{y\ast(\Psi f)(x) }_{(3)}
\]
\textbf{Term (1): $x \ast (\Psi f)(y)$}


\begin{align*}
&= (x_1 \otimes b_1) \ast \left(x_2 \otimes f(b_2) \right) \\
&= [x_1, x_2 ] \otimes (b_1  f(b_2)) - [x_2, x_1] \otimes ( f(b_2)b_1) 
\end{align*}
\textbf{Term (2): $-(\Psi f)(x \ast y)$}

\begin{align*}
&=-f[ (x_1 \otimes b_1) \ast (x_2 \otimes b_2)] \\
&=- [x_1, x_2 ] \otimes f(b_1b_2) + [x_2, x_1] \otimes f(b_2b_1) 
\end{align*}

\textbf{Term (3): $-y\ast (\Psi f)(x) $}
\begin{align*}
&=-(x_2 \otimes b_2) \ast (x_1 \otimes f(b_1)) \\
&= -[x_2, x_1 ] \otimes b_2f(b_1) + [x_1, x_2] \otimes f(b_1)b_2 
\end{align*}
Now, let's collect all terms from (1)-(3) and group them.

\[
\begin{aligned}
\delta_{Lie}(\Psi f)(x,y)&=[x_1,x_2]\otimes \Big( b_1  f(b_2)-f(b_1b_2)+ f(b_1)b_2\Big)-[x_2,x_1]\otimes \Big( b_2f(b_1)  -f(b_2b_1) +f(b_2)b_1\Big)\\
&=[x_1,x_2]\otimes (\delta_{\DL}f)(b_1,b_2)-[x_2,x_1]\otimes (\delta_{\DL}f)(b_2,b_1)\\
&=\Psi(\delta_{\DL}f)(x,y)
\end{aligned}
\]
For $n=3$ and $x =x_1  \otimes b_1$, $y = x_2 \otimes b_2$, $z = x_3 \otimes b_3$, we have: $\delta_{lie}(\Psi f)(x,y,z) =$

\[
 \underbrace{x \ast (\Psi f)(y,z)}_{(1)} -\underbrace{y \ast (\Psi f)(x,z)}_{(2)}+\underbrace{z \ast (\Psi f)(x,y)}_{(3)}- \underbrace{(\Psi f)(x \ast y, z)}_{(4)} + \underbrace{(\Psi f)(x\ast z, y )}_{(5)}-\underbrace{(\Psi f)(y\ast z, x )}_{(6)} 
\]
\textbf{Term (1): $x \ast (\Psi f)(y,z)$}

\begin{align*}
&= (x_1 \otimes b_1) \ast \Big([x_2 , x_3] \otimes f(b_2,b_3) - [x_3, x_2] \otimes f(b_3,b_2)\Big) \\
&= [x_1, [x_2 , x_3]] \otimes (b_1  f(b_2,b_3)) - [[x_2, x_3], x_1] \otimes ( f(b_2,b_3)b_1) \\
&\quad - [x_1, [x_3, x_2]] \otimes (b_1  f(b_3,b_2)) + [[x_3, x_2], x_1] \otimes (f(b_3,b_2)b_1)\\
&=[[x_1, x_2], x_3]] \otimes (b_1  f(b_2,b_3))-[[x_1, x_3], x_2]] \otimes (b_1  f(b_2,b_3)) - [[x_2, x_3],x_1] \otimes (f(b_2,b_3)b_1) \\
&\quad - [[x_1, x_3], x_2]] \otimes (b_1  f(b_3,b_2))+[[x_1, x_2], x_3]] \otimes (b_1  f(b_3,b_2)) + [[x_3, x_2], x_1] \otimes (f(b_3,b_2)b_1)\\
\end{align*}
\textbf{Term (2): $-y \ast (\Psi f)(x,z)$}

\begin{align*}
&= -(x_2 \otimes b_2) \ast \Big([x_1 , x_3] \otimes f(b_1,b_3) + [x_3, x_1] \otimes f(b_3,b_1)\Big) \\
&= -[x_2, [x_1, x_3]] \otimes (b_2  f(b_1,b_3)) + [[x_1, x_3], x_2] \otimes ( f(b_1,b_3)b_2) \\
&\quad + [x_2, [x_3, x_1]] \otimes (b_2  f(b_3,b_1)) - [[x_3, x_1], x_2] \otimes (f(b_3,b_1)b_2)\\
&= -[[x_2, x_1], x_3]] \otimes (b_2  f(b_1,b_3))+[[x_2, x_3], x_1]] \otimes (b_2  f(b_1,b_3)) + [[x_1, x_3], x_2] \otimes ( f(b_1,b_3)b_2) \\
&\quad + [[x_2, x_3], x_1]] \otimes (b_2  f(b_3,b_1))- [[x_2, x_1], x_3]] \otimes (b_2  f(b_3,b_1)) - [[x_3, x_1], x_2] \otimes (f(b_3,b_1)b_2)\\
\end{align*}
\textbf{Term (3): $z \ast (\Psi f)(x,y)$}

\begin{align*}
&= (x_3 \otimes b_3) \ast \Big([x_1 , x_2] \otimes f(b_1,b_2) -[x_2, x_1] \otimes f(b_2,b_1)\Big) \\
&= [x_3, [x_1, x_2]] \otimes (b_3  f(b_1,b_2)) - [[x_1, x_2], x_3] \otimes ( f(b_1,b_2)b_3) \\
&\quad - [x_3, [x_2, x_1]] \otimes (b_3  f(b_2,b_1)) + [[x_2, x_1], x_3] \otimes (f(b_2,b_1)b_3)\\
&= [[x_3, x_1], x_2]] \otimes (b_3  f(b_1,b_2))-[[x_3, x_2], x_1]] \otimes (b_3  f(b_1,b_2)) - [[x_1, x_2], x_3] \otimes ( f(b_1,b_2)b_3) \\
&\quad - [[x_3, x_2], x_1]] \otimes (b_3  f(b_2,b_1))+[[x_3, x_1], x_2]] \otimes (b_3  f(b_2,b_1)) + [[x_2, x_1], x_3] \otimes (f(b_2,b_1)b_3)\\
\end{align*}
\textbf{Term (4): $-(\Psi f)(x \ast y, z)$}

\begin{align*}
&= -(\Psi f)\Big([x_1, x_2] \otimes (b_1 b_2) - [x_2, x_1] \otimes (b_2 b_1), x_3 \otimes b_3\Big) \\
&= - [[x_1, x_2], x_3] \otimes f( b_1 b_2, b_3) + [x_3, [x_1, x_2]] \otimes f(b_3,b_1 b_2) \\
&\quad + [[x_2, x_1], x_3] \otimes f(b_2 b_1, b_3) -[x_3, [x_2, x_1]] \otimes f( b_3,b_2  b_1)\\
&= - [[x_1, x_2], x_3] \otimes f( b_1 b_2, b_3) + [[x_3, x_1], x_2]] \otimes f(b_3,b_1 b_2)-[[x_3, x_2], x_1]] \otimes f(b_3,b_1 b_2) \\
&\quad + [[x_2, x_1], x_3] \otimes f(b_2 b_1, b_3) -[[x_3, x_2], x_1]] \otimes f( b_3,b_2  b_1)+[[x_3, x_1], x_2]] \otimes f( b_3,b_2  b_1)
\end{align*}
\textbf{Term (5): $+(\Psi f)(x \ast z, y)$}

\begin{align*}
&= (\Psi f)\Big([x_1, x_3] \otimes (b_1 b_3) - [x_3, x_1] \otimes (b_3 b_1), x_2 \otimes b_2\Big) \\
&=  [[x_1, x_3], x_2] \otimes f( b_1 b_3, b_2) - [x_2, [x_1, x_3]] \otimes f(b_2,b_1 b_3) \\
&\quad - [[x_3, x_1], x_2] \otimes f(b_3 b_1, b_2) +[x_2, [x_3, x_1]] \otimes f( b_2,b_3  b_1)\\
&=  [[x_1, x_3], x_2] \otimes f( b_1 b_3, b_2) - [[x_2, x_1], x_3]] \otimes f(b_2,b_1 b_3)+[[x_2, x_3], x_1]] \otimes f(b_2,b_1 b_3) \\
&\quad - [[x_3, x_1], x_2] \otimes f(b_3 b_1, b_2) +[[x_2, x_3], x_1]] \otimes f( b_2,b_3  b_1)-[[x_2, x_1], x_3]] \otimes f( b_2,b_3  b_1)
\end{align*}
\textbf{Term (6): $-(\Psi f)(y \ast z, x)$}

\begin{align*}
&= -(\Psi f)\Big([x_2, x_3] \otimes (b_2 b_3) - [x_3, x_2] \otimes (b_3 b_2), x_1 \otimes b_1\Big) \\
&=  -[[x_2, x_3], x_1] \otimes f( b_2 b_3, b_1) + [x_1, [x_2, x_3]] \otimes f(b_1,b_2 b_3) \\
&\quad + [[x_3, x_2], x_1] \otimes f(b_3 b_2, b_1) -[x_1, [x_3, x_2]] \otimes f( b_1,b_3  b_2)\\
&=  -[[x_2, x_3], x_1] \otimes f( b_2 b_3, b_1) + [[x_1, x_2], x_3]] \otimes f(b_1,b_2 b_3)-[[x_1, x_3], x_2]] \otimes f(b_1,b_2 b_3) \\
&\quad + [[x_3, x_2], x_1] \otimes f(b_3 b_2, b_1) -[[x_1, x_3], x_2]] \otimes f( b_1,b_3  b_2)+[[x_1, x_2], x_3]] \otimes f( b_1,b_3  b_2)
\end{align*}


Now, let's collect all terms from (1)-(6) and group them.

\textbf{Terms Contributing to $[[x_1,x_2],x_3]$:}

\begin{itemize}
\item From (1): $b_1f(b_2,b_3)+b_1  f(b_3,b_2)$
\item From (3): $-f(b_1, b_2)b_3$
\item From (4): $-f(b_1b_2, b_3)$ 
\item From (6): $+f(b_1,b_2b_3)+f(b_1,b_3b_2)$
\end{itemize}

Combined:
\begin{align*}
(\delta_{\DL} f)(b_1, b_2, b_3) &= b_1 f(b_2,b_3)+b_1 f(b_3,b_2) - f(b_1  b_2, b_3) \\
&\quad + f(b_1, b_2  b_3) + f(b_1, b_3  b_2) - f(b_1,b_2)  b_3
\end{align*}
\textbf{Terms Contributing to $-[[x_1,x_3],x_2]$:}

\begin{itemize}
\item From (1): $b_1f(b_2,b_3)+b_1f(b_3,b_2)$
\item From (2): $-f(b_1, b_3)b_2$
\item From (5): $-f(b_1b_3, b_2)$
\item From (4): $+f(b_1,b_2b_3)+f(b_1,b_3b_2)$
\end{itemize}
Combined:
\begin{align*}
(\delta_{\DL} f)(b_1, b_3, b_2) &= b_1f(b_2,b_3)+b_1 f(b_3,b_2) - f(b_1  b_3, b_2) \\
&\quad + f(b_1, b_2  b_3) + f(b_1, b_3  b_2) - f(b_1,b_3)  b_2
\end{align*}
\textbf{Terms Contributing to $[[x_2,x_3],x_1]$:}
\begin{itemize}
\item From (1): $-f(b_2,  b_3)b_1$ 
\item From (2):  $  b_2f(b_1,b_3)+b_2f(b_3,b_1)$
\item From (5): $+f(b_2, b_3  b_1)+f(b_2,b_1b_3)$
\item From (6): $-f(b_2b_3,b_1)$
\end{itemize}
Combined:
\begin{align*}
(\delta_{\DL} f)(b_2, b_3, b_1) &= b_2f(b_3,b_1)+b_2 f(b_1,b_3) - f(b_2  b_3, b_1) \\
&\quad + f(b_2, b_3  b_1) + f(b_2, b_1  b_3) - f(b_2,b_3)  b_1
\end{align*}
\textbf{Terms Contributing to $-[[x_2,x_1],x_3]$:}
\begin{itemize}
\item From (2): $b_2f(b_1,  b_3)+b_2f(b_3,b_1)$ 
\item From (3):  $ - f(b_2,b_1)b_3$
\item From (4): $-f(b_2b_1, b_3 )$
\item From (6): $+f(b_2,b_3b_1)+f(b_2,b_1b_3)$
\end{itemize}
Combined:
\begin{align*}
(\delta_{\DL} f)(b_2, b_1, b_3) &= b_2f(b_3,b_1)+b_2 f(b_1,b_3) - f(b_2  b_1, b_3) \\
&\quad + f(b_2, b_3  b_1) + f(b_2, b_1  b_3) - f(b_2,b_1)  b_3
\end{align*}
\textbf{Terms Contributing to $[[x_3,x_1],x_2]$:}

\begin{itemize}
\item From (2): $-f(b_3,b_1)b_2$
\item From (3): $+b_3f(b_1, b_2)+b_3f(b_2,b_1)$
\item From (4): $+f(b_3,b_1b_2)+f(b_3,b_2b_1)$ 
\item From (6): $-f(b_3b_1,b_2)$
\end{itemize}

Combined:
\begin{align*}
(\delta_{\DL} f)(b_3, b_1, b_2) &= b_3 f(b_1,b_2)+b_3 f(b_2,b_1) - f(b_3  b_1, b_2) \\
&\quad + f(b_3, b_1  b_2) + f(b_3, b_2  b_1) - f(b_3,b_1)  b_2
\end{align*}
\textbf{Terms Contributing to $-[[x_3,x_2],x_1]$:}

\begin{itemize}
\item From (1): $-f(b_3,b_2)b_1$
\item From (3): $b_3f(b_1, b_2)+b_3f(b_2,b_1)$
\item From (4): $+f(b_3,b_1b_2)+f(b_3,b_2b_1)$ 
\item From (6): $-f(b_3b_2,b_2)$
\end{itemize}

Combined:
\begin{align*}
(\delta_{\DL} f)(b_3, b_2, b_1) &= b_3 f(b_2,b_1)+b_3 f(b_1,b_2) - f(b_3  b_2, b_1) \\
&\quad + f(b_3, b_2  b_1) + f(b_3, b_1  b_2) - f(b_3,b_2)  b_1
\end{align*}
Therefore
\[
\begin{aligned}
\delta_{Lie}(\Psi f)(x,y,z) &=[[x_1,x_2],x_3]\otimes (\delta_{\DL}f)(b_1,b_2,b_3)-[[x_1,x_3],x_2]\otimes (\delta_{\DL}f)(b_1,b_3,b_2)\\
&+[[x_2,x_3],x_1]\otimes (\delta_{\DL}f)(b_2,b_3,b_1)-[[x_2,x_1],x_3]\otimes (\delta_{\DL}f)(b_2,b_1,b_3)\\
&+[[x_3,x_1],x_2]\otimes (\delta_{\DL}f)(b_3,b_1,b_2)-[[x_3,x_2],x_1]\otimes (\delta_{\DL}f)(b_3,b_2,b_1)\\
&=\sum_{\sigma\in S_3}sgn(\sigma)[[x_{\sigma(1)},x_{\sigma(2)}],x_{\sigma(3)}]\otimes \delta_{\Zinb}(b_{\sigma(1)},b_{\sigma(2)},b_{\sigma(3)})=\Psi(\delta_{Lie} f)(x,y,z)
\end{aligned}
\]









\textbf{General case:}
For a cochain $\Psi f$ of degree $n$ can be writting as: 
\[
\Psi(f)(x_1 \otimes b_1, \dots, x_n \otimes b_n) = \sum_{\sigma \in S_n} \operatorname{sgn}(\sigma) L_{x_\sigma} \otimes f(b_{\sigma})
\]
where \(L_{x_\sigma} = [[[x_{\sigma(1)}, x_{\sigma(2)}], x_{\sigma(3)}], \cdots, x_{\sigma(n)}]\) is the left-normed bracket and $b_{\sigma}=(b_{\sigma(1)},\cdots,b_{\sigma(n)})$.


For a cochain $\Psi f$ of degree $n$, the Lie differential \(\delta_{Lie}(\Psi f)\) is given by:
\[
\begin{aligned}
(\delta_{\mathrm{Lie}} \Psi f)(y_1, \dots, y_{n+1}) = &\ \sum_{i=1}^{n+1} (-1)^{i+1} y_i \ast (\Psi f)(y_1, \dots, \widehat{y_i}, \dots, y_{n+1}) \\
& +  \sum_{1\leq i < j} (-1)^{i+j} (\Psi f)(y_i\ast y_j, y_1, \ldots, \hat{y_i}, \ldots, \hat{y_j}, \ldots, y_{n+1})
\end{aligned}
\]


Let \(y_i = x_i \otimes b_i\) for \(i = 1, \dots, n+1\). We denote the first sum as Term I (the action terms) and the second as Term II (the bracket terms).

Our goal is to show that for all \(y_1, \dots, y_{n+1}\),
\[
\delta_{\mathrm{Lie}} (\Psi f)(y_1, \dots, y_{n+1}) = \Psi(\delta_{\DL} f)(y_1, \dots, y_{n+1}).
\]

\textbf{Term I: Action Terms}
\[
I = \sum_{i=1}^{n+1} (-1)^{i+1} y_i \ast \Psi(f)(y_1, \dots, \widehat{y_i}, \dots, y_{n+1})
\]

Let us analyze one summand. For fixed \(i\):
\[
\Psi(f)(y_1, \dots, \widehat{y_i}, \dots, y_{n+1}) = \sum_{\tau \in S_n} \operatorname{sgn}(\tau) L_{x_\tau} \otimes f(b_{\tau})
\]
where \(\tau\) permutes \(\{1, \dots, n+1\} \setminus \{i\}\), and \(L_{x_\tau}\) is the left-normed bracket of the corresponding \(x\)-elements.

Now, for a fixed \(\tau\), consider the action:
\[
y_i \ast (L_{x_\tau} \otimes f(\mathbf{b}_\tau)).
\]
Using the module structure on \(\mathfrak{g} \otimes M\), this equals:
\[
[x_i, L_{x_\tau}] \otimes (b_i \cdot f(\mathbf{b}_\tau)) \;-\; [L_{x_\tau}, x_i] \otimes (f(\mathbf{b}_\tau) \cdot b_i). \tag{1}
\]

We analyze these two parts separately.

\textbf{Part I.A: Contribution from \([x_i, L_{x_\tau}] \otimes (b_i \cdot f(\mathbf{b}_\tau))\)}

We use the following identity (see Lemma \ref{lem}), which is a consequence of the Leibniz identity applied iteratively to left-normed brackets. For any \(z, w_1, \dots, w_n \in \mathfrak{g}\),
\[
[z, [w_1, \dots, w_n]] = \sum_{k=0}^{n-1} \sum_{(\alpha, \beta) \in \mathsf{Sh}(n-k, k)} (-1)^k \, [z, w_{\beta(k)}, \dots, w_{\beta(1)}, w_{\alpha(1)}, \dots, w_{\alpha(n-k)}], \tag{2}
\]
where the inner bracket on the right is also left-normed, and the sign comes from the shuffle.

Apply this with \(z = x_i\) and \((w_1, \dots, w_n) = (x_{\tau(1)}, \dots, x_{\tau(n)})\). For each \(k\) and each \((n-k, k)\)-shuffle \((\alpha, \beta)\), we get a term:
\[
\operatorname{sgn}(\tau) \cdot (-1)^k \cdot [x_i, x_{\tau(\beta(k))}, \dots, x_{\tau(\beta(1))}, x_{\tau(\alpha(1))}, \dots, x_{\tau(\alpha(n-k))}] \otimes (b_i \cdot f(b_{\tau(1)}, \dots, b_{\tau(n)})).
\]
Define a new permutation \(\sigma \in S_{n+1}\) by:
\[
\sigma = (i, \tau(\beta(k)), \dots, \tau(\beta(1)), \tau(\alpha(1)), \dots, \tau(\alpha(n-k))).
\]


Furthermore, the coefficient \(b_i \cdot f(b_{\tau(1)}, \dots, b_{\tau(n)})\) corresponds to \(b_{\sigma(1)} \cdot f(b_{\sigma(2)}, \dots, b_{\sigma(n+1)})\), since \(\sigma(1) = i\).

Thus, for a fixed \(\sigma \in S_{n+1}\), the total contribution to the coefficient of \(L_{x_\sigma} \otimes (-)\) from Part I.A comes from all triples \((i, \tau, (\alpha, \beta))\) which produce this specific \(\sigma\), which is denoted by $\overline{S_{n}}$ and $\overline{\operatorname{sgn}}\tau=\operatorname{sgn}\tau(-1)^k$. The sum of the signs and the scalar part yields a contribution of:
\[
L_{x_\sigma} \otimes \left( \operatorname{sgn}(\sigma) b_{\sigma(1)}.  \sum_{\tau \in \overline{S_{n}}}\overline{\operatorname{sgn}}(\tau)f(b_{\tau(2)}, \dots, b_{\tau(n+1)}) \right).
\]
The alternating sign \((-1)^{i+1}\) from Term I combines with the sign from the bracket identity to produce the correct overall sign \(\operatorname{sgn}(\sigma)\) for the term in the sum defining \(\Psi(\delta_{\DL} f)\). A careful but standard combinatorial argument confirms this.

\textbf{Part I.B: Contribution from \(- [L_{x_\tau}, x_i] \otimes (f(\mathbf{b}_\tau) \cdot b_i)\)}


Define a permutation \(\sigma \in S_{n+1}\) by:
\[
\sigma = (\tau(1), \dots, \tau(n), i).
\]
Then \(L_{x_\sigma} = L_{(x_\tau, x_i)} = [L_{x_\tau}, x_i]\). The sign of \(\sigma\) is \(\operatorname{sgn}(\sigma) = \operatorname{sgn}(\tau)\) (since \(i\) is fixed in the last position). The scalar part is \(f(\mathbf{b}_\tau) \cdot b_i = f(b_{\sigma(1)}, \dots, b_{\sigma(n)}) \cdot b_{\sigma(n+1)}\).

Therefore, the term \(- [L_{x_\tau}, x_i] \otimes (f(\mathbf{b}_\tau) \cdot b_i)\) contributes:
\[
- \operatorname{sgn}(\tau) \, L_{x_\sigma} \otimes (f(b_{\sigma(1)}, \dots, b_{\sigma(n)}) \cdot b_{\sigma(n+1)}).
\]
The sign from Term I is \((-1)^{i+1}\). Since \(i = \sigma(n+1)\), we have \((-1)^{i+1} = (-1)^{\sigma(n+1)+1}\). Combining this with \(-\operatorname{sgn}(\tau) = -\operatorname{sgn}(\sigma)\), the total sign is \((-1)^{\sigma(n+1)} \operatorname{sgn}(\sigma)\). For the term to appear in \(\Psi(\delta_{\DL} f)\) with sign \(\operatorname{sgn}(\sigma)\), we need a factor of \((-1)^{n+1}\). Indeed, \((-1)^{\sigma(n+1)}\) is not simply \((-1)^{n+1}\), but when summed over all \(\tau\) that lead to the same \(\sigma\), the overall contribution consolidates to:
\[
L_{x_\sigma} \otimes \left( (-1)^{n+1} f(b_{\sigma(1)}, \dots, b_{\sigma(n)}) \cdot b_{\sigma(n+1)} \right).
\]

\textbf{Summary of Term I:}
For a fixed \(\sigma \in S_{n+1}\), Term I contributes the following to the coefficient of \(L_{x_\sigma} \otimes (-)\):
\[
L_{x_\sigma} \otimes \left(\operatorname{sgn}(\sigma) \Big( b_{\sigma(1)}.  \sum_{\tau \in \overline{S_{n}}}\overline{\operatorname{sgn}}(\tau)f(b_{\tau(2)}, \dots, b_{\tau(n+1)})+ (-1)^{n+1} f(b_{\sigma(1)}, \dots, b_{\sigma(n)}) \cdot b_{\sigma(n+1)} \Big) \right). \tag{I}
\]




\textbf{ Term II - Bracket Terms:}

\[
II = \sum_{1 \leq i < j \leq n+1} (-1)^{i+j} \, (\Psi f)(y_i \ast y_j, y_1, \dots, \widehat{y_i}, \dots, \widehat{y_j}, \dots, y_{n+1}).
\]
Fix a pair \(i < j\). Then:
\[
y_i \ast y_j = [x_i, x_j] \otimes (b_i b_j) - [x_j, x_i] \otimes (b_j b_i).
\]
Thus,
\[
\begin{aligned}
&(\Psi f)(y_i \ast y_j, y_1, \dots, \widehat{y_i}, \dots, \widehat{y_j}, \dots, y_{n+1}) = \\
&\quad (\Psi f)([x_i, x_j] \otimes (b_i b_j), \dots) \;-\; (\Psi f)([x_j, x_i] \otimes (b_j b_i), \dots).
\end{aligned}
\]
We analyze the first term, \((\Psi f)([x_i, x_j] \otimes (b_i b_j), \dots)\). Let \(\rho\) be a permutation of the \(n-1\) indices \(\{1, \dots, n+1\} \setminus \{i, j\}\). The cochain \(\Psi(f)\) evaluated on the \(n\) arguments \(([x_i, x_j] \otimes (b_i b_j), y_{\rho(1)}, \dots, y_{\rho(n-1)})\) is:
\[
\sum_{\pi \in S_n} \operatorname{sgn}(\pi) \, L_{z_\pi} \otimes f(c_{\pi(1)}, \dots, c_{\pi(n)}),
\]
where the list \((z_1, c_1), \dots, (z_n, c_n)\) is \(([x_i, x_j], b_i b_j), (x_{\rho(1)}, b_{\rho(1)}), \dots, (x_{\rho(n-1)}, b_{\rho(n-1)})\).

The key is that the bracket \(L_{z_\pi}\) is linear. If \(\pi(1) = 1\), meaning the first argument is \(([x_i, x_j], b_i b_j)\), then
\[
L_{z_\pi} = [[x_i, x_j], x_{\rho(\pi(2)-1)}, \dots, x_{\rho(\pi(n)-1)}].
\]

















 In case  \(\pi(1) \neq 1\), we use the identity for the left-normed bracket when the middle element is a bracket:
\[
[ x_{k_1}, \dots,[x_i, x_j],\dots, x_{k_{n-1}}] = [ x_{k_1}, \dots, x_i,x_j,\dots, x_{k_{n-1}}] - [ x_{k_1}, \dots, x_j,x_i,\dots, x_{k_{n-1}}].
\]
Thus, the term where the first argument is the bracket contributes:
\[
\operatorname{sgn}(\pi) \left( [ x_{\rho(\pi(2)-1)}, \dots,x_i, x_j, \dots, x_{\rho(\pi(n)-1)}] - [ x_{\rho(\pi(2)-1)}, \dots, x_j, x_i,\dots, x_{\rho(\pi(n)-1)}] \right)
\]

Now, define a permutation \(\sigma \in S_{n+1}\) as follows. For the first term, let:
\[
\sigma = ( \rho(\pi(2)-1), \dots,i, j,\dots, \rho(\pi(n)-1)).
\]
Then \(L_{x_\sigma} = [ x_{\rho(\pi(2)-1)}, \dots, x_i, x_j,\dots,x_{\rho(\pi(n)-1)}]\).

The sign is \(\operatorname{sgn}(\sigma) = \operatorname{sgn}(\pi) \cdot \varepsilon_1\), where \(\varepsilon_1\) is the sign for placing \(i, j\) first. The scalar part is \(
f(b_{\sigma(1)},\cdots, b_{\sigma(i)}b_{\sigma(j)}, \dots, b_{\sigma(n+1)})\).

Similarly, the second term corresponds to \(\sigma' = ( \rho(\pi(2)-1), \dots,j, i,\dots, \rho(\pi(n)-1))\), with \(\operatorname{sgn}(\sigma') = -\operatorname{sgn}(\sigma)\). The scalar part is \(
f(b_{\sigma(1)},\cdots, b_{\sigma(j)}b_{\sigma(i)}, \dots, b_{\sigma(n+1)})\).

The sign \((-1)^{i+j}\) from Term II combines with the signs from the permutation and the bracket expansion to yield the correct overall sign \(\operatorname{sgn}(\sigma)\).


A systematic analysis shows that for a fixed \(\sigma \in S_{n+1}\), the contributions from Term II to the coefficient of \(L_{x_\sigma} \otimes (-)\) are precisely:
\[
L_{x_\sigma} \otimes \left(\operatorname{sgn}\sigma \Big( \sum_{k=1}^{n} (-1)^k f(b_{\sigma(1)}, \dots, b_{\sigma(k)} b_{\sigma(k+1)}, \dots, b_{\sigma(n+1)}) \;+\; \sum_{k=2}^{n} (-1)^k f(b_{\sigma(1)}, \dots, b_{\sigma(k+1)} b_{\sigma(k)}, \dots, b_{\sigma(n+1)}) \Big)\right). \tag{II}
\]
The factor \((-1)^k\) accounts for the alternating sign in the Lie coboundary and the position of the product in the arguments of \(f\).

\textbf{ Combining Terms I and II:} For a fixed permutation \(\sigma \in S_{n+1}\), we now combine the contributions (I) and (II). The total coefficient of \(\operatorname{sgn}\sigma. L_{x_\sigma} \otimes (-)\) in \((\delta_{\mathrm{Lie}} \Psi f)(y_1, \dots, y_{n+1})\) is:
\[
\begin{aligned}
&  b_{\sigma(1)}.  \sum_{\tau \in \overline{S_{n+1}}- \{\sigma(1)\}}\operatorname{sgn}(\tau)f(b_{\tau(2)}, \dots, b_{\tau(n+1)}) \quad &\text{(from I)} \\
&+ \sum_{k=1}^{n} (-1)^k f(b_{\sigma(1)}, \dots, b_{\sigma(k)} b_{\sigma(k+1)}, \dots, b_{\sigma(n+1)}) \quad &\text{(from II, first sum)} \\
&+ \sum_{k=1}^{n} (-1)^k f(b_{\sigma(1)}, \dots, b_{\sigma(k+1)} b_{\sigma(k)}, \dots, b_{\sigma(n+1)}) \quad &\text{(from II, second sum)} \\
&+ (-1)^{n+1} f(b_{\sigma(1)}, \dots, b_{\sigma(n)}) \cdot b_{\sigma(n+1)}. \quad &\text{(from I)}
\end{aligned}
\]
But this is exactly \((\delta_{\DL} f)(b_{\sigma(1)}, \dots, b_{\sigma(n+1)})\), the dual Leibniz coboundary of \(f\) evaluated on \((b_{\sigma(1)}, \dots, b_{\sigma(n+1)})\).

Therefore, we have:
\[
\delta_{\mathrm{Lie}} (\Psi f)(y_1, \dots, y_{n+1}) = \sum_{\sigma \in S_{n+1}} \operatorname{sgn}(\sigma) \, L_{x_\sigma} \otimes (\delta_{\DL} f)(b_{\sigma(1)}, \dots, b_{\sigma(n+1)}) = \Psi(\delta_{\DL} f)(y_1, \dots, y_{n+1}).
\]

This completes the proof that \(\Psi\) is a cochain map.

 
\end{proof}

\begin{thebibliography}{99}
\bibitem[I.S]{I.S} I. Alekseev and S.O. Ivanov, \emph{Higher Jacobi identities}, https://doi.org/10.48550/arXiv.1604.05281.
\end{thebibliography}
\end{document}
